% Fragmentation example (interleaved, all-manual parameterization). Two layers:
%  - top: each step's difficulty for AI (easy=teal, hard=red), alternating;
%  - bottom: the optimal execution of each step (all manual=gray).
% Standalone augmentation does not pay (management cost exceeds manual cost) and
% no two easy steps are adjacent, so no chain forms: every step is manual and the
% workflow uses no AI at all. Colors follow Figures 1-2.
\centering
\begin{adjustbox}{max width=\linewidth}
\begin{tikzpicture}[>=latex,
    box/.style={rectangle, rounded corners, draw, minimum width=1.8cm, minimum height=0.8cm, align=center, font=\small},
    easy/.style={fill=teal!40},
    hard/.style={fill=red!40},
    manual/.style={fill=gray!10},
    idx/.style={draw=none, font=\scriptsize},
    rowlab/.style={draw=none, font=\footnotesize\itshape}]

  % ---- top layer: step difficulty ----
  \node[box, easy] (d1) {Easy};
  \node[box, hard, right=0.7cm of d1] (d2) {Hard};
  \node[box, easy, right=0.7cm of d2] (d3) {Easy};
  \node[box, hard, right=0.7cm of d3] (d4) {Hard};

  \node[idx, above=1pt of d1] {Step 1};
  \node[idx, above=1pt of d2] {Step 2};
  \node[idx, above=1pt of d3] {Step 3};
  \node[idx, above=1pt of d4] {Step 4};

  % ---- bottom layer: optimal execution (tasks) ----
  \node[box, manual, below=1.4cm of d1] (e1) {Manual};
  \node[box, manual, below=1.4cm of d2] (e2) {Manual};
  \node[box, manual, below=1.4cm of d3] (e3) {Manual};
  \node[box, manual, below=1.4cm of d4] (e4) {Manual};

  \foreach \a/\b in {d1/e1,d2/e2,d3/e3,d4/e4} \draw[->,thick] (\a)--(\b);

  \node[rowlab, left=0.5cm of d1] {Difficulty:};
  \node[rowlab, left=0.5cm of e1] {Execution:};
\end{tikzpicture}
\end{adjustbox}
